\documentclass[11pt]{article}
\usepackage[margin=1in]{geometry}
\usepackage{amsmath,amssymb,amsthm,hyperref}
\newtheorem{theorem}{Theorem}
\newtheorem{lemma}[theorem]{Lemma}
\newtheorem{corollary}[theorem]{Corollary}
\newcommand{\E}{\mathbb E}
\newcommand{\R}{\mathbb R}
\DeclareMathOperator{\erfc}{erfc}
\title{Complete monotonicity of a squared Mittag--Leffler modulus}
\author{Principia Math In Progress}
\date{11 September 2026}
\begin{document}
\maketitle
\begin{abstract}
For every $0<\alpha\leq1/2$, the function
$|E_{2\alpha,1}(-t)|^2+t|E_{2\alpha,1+\alpha}(-t)|^2$
is strictly completely monotone and strictly log-convex on $(0,\infty)$.
This proves Conjecture 3 of Chorfi, Et-tahri, Maniar, and Yamamoto.
The proof combines an inverse-stable scale representation with an explicit
positive Laplace representation for differences of scaled half-Gaussians.
It also gives the proposed logarithmic convexity estimate for the
corresponding spectral solutions.
\end{abstract}

\section*{Task 1: Resolve MathDB problem 370772}
For $\rho,\delta>0$, write
\[
 E_{\rho,\delta}(z)=\sum_{n=0}^\infty\frac{z^n}{\Gamma(\delta+\rho n)},
 \qquad E_\alpha=E_{\alpha,1}.
\]
Conjecture 3, equation (5.2), page 23 of \cite{source}, asks whether
\begin{equation}\label{eq:target}
 f_\alpha(t)=|E_{2\alpha,1}(-t)|^2
       +t|E_{2\alpha,1+\alpha}(-t)|^2,
 \qquad 0<\alpha\leq\tfrac12,
\end{equation}
is completely monotone, or at least log-convex, for $t>0$.
The authors still refer to the journal version's Conjecture 3 as unresolved
in their March 2026 follow-up \cite[page 13]{followup}.

\begin{theorem}\label{thm:main}
For every $0<\alpha\leq1/2$, there is a nondegenerate probability measure
$\nu_\alpha$ on $(0,\infty)$ such that
\begin{equation}\label{eq:laplace}
 f_\alpha(t)=\int_0^\infty e^{-t\lambda}\,\nu_\alpha(d\lambda)
 \quad(t\geq0).
\end{equation}
Consequently, $(-1)^n f_\alpha^{(n)}(t)>0$ for every integer $n\geq0$
and $t>0$, and $(\log f_\alpha)''(t)>0$ there.
\end{theorem}

\subsection*{A half-Gaussian scale representation}
Let $H$ have density $\pi^{-1/2}e^{-x^2/4}$ on $x>0$.
For $0<\beta<1$, let $S_\beta$ be a positive stable random variable,
normalized by $\E e^{-sS_\beta}=e^{-s^\beta}$ for $s\geq0$;
set $S_1=1$.
These standard laws and their Mittag--Leffler transforms are discussed in
\cite[Examples 3.10--3.11]{janson}. We give the moment calculation needed here.
For $q>0$, Tonelli's theorem and $u=s^\beta$ give
\begin{equation}\label{eq:negative-moment}
 \E S_\beta^{-q}
 =\frac1{\Gamma(q)}\int_0^\infty s^{q-1}e^{-s^\beta}\,ds
 =\frac{\Gamma(1+q/\beta)}{\Gamma(1+q)}.
\end{equation}
The formula also holds for $\beta=1$.

\begin{lemma}\label{lem:scale}
Fix $0<\alpha\leq1/2$, let $R=S_{2\alpha}^{-\alpha}$, and take $H$
independent of $R$. Then $X=RH$ satisfies
\begin{equation}\label{eq:mgf}
 \E e^{zX}=E_\alpha(z)\qquad(z\in\mathbb C).
\end{equation}
\end{lemma}
\begin{proof}
For every integer $n\geq0$, the gamma duplication formula gives
\[
 \E H^n=\frac{2^n\Gamma((n+1)/2)}{\sqrt\pi}
       =\frac{n!}{\Gamma(1+n/2)},\qquad
 \E R^n=\frac{\Gamma(1+n/2)}{\Gamma(1+\alpha n)}.
\]
The second identity follows from \eqref{eq:negative-moment} for $n>0$
and is immediate for $n=0$. Independence yields
$\E X^n=n!/\Gamma(1+\alpha n)$.
For $c\geq0$, monotone convergence therefore gives
\[
 \E e^{cX}=\sum_{n=0}^\infty\frac{c^n}{\Gamma(1+\alpha n)}<\infty;
\]
convergence follows from Stirling's formula. For complex $z$, absolute
integrability follows by the same calculation with $c=|z|$.
Termwise expectation proves \eqref{eq:mgf}.
\end{proof}

\subsection*{A positive Laplace kernel}
\begin{lemma}\label{lem:kernel}
Let $H,H'$ be independent copies of the half-Gaussian above. For any
$a,b>0$, set $A=a^2+b^2$. Then
\begin{align}
 K_{a,b}(t)&:=\E\cos\bigl(\sqrt t(aH-bH')\bigr)\notag\\
 &=\frac2\pi\int_0^\infty\left[
 \frac{a\exp\{-a^2At/(A+b^2v^2)\}}
 {(1+v^2)\sqrt{A+b^2v^2}}
 +\frac{b\exp\{-b^2At/(A+a^2v^2)\}}
 {(1+v^2)\sqrt{A+a^2v^2}}\right]dv.\label{eq:kernel}
\end{align}
This is the Laplace transform of a nondegenerate probability measure
$\nu_{a,b}$ on $(0,\infty)$.
\end{lemma}
\begin{proof}
Let $h_{a,b}$ be the density of $aH-bH'$. For $z\geq0$, completing the
square in the convolution integral gives
\begin{align*}
 h_{a,b}(z)
 &=\frac1{ab\pi}\int_0^\infty
    e^{-(z+y)^2/(4a^2)-y^2/(4b^2)}\,dy\\
 &=\frac{e^{-z^2/(4A)}}{\sqrt{\pi A}}
   \erfc\left(\frac{bz}{2a\sqrt A}\right),
 \qquad h_{a,b}(-z)=h_{b,a}(z).
\end{align*}
Thus the symmetrized probability density is
\begin{equation}\label{eq:density}
 k_{a,b}(z):=\frac{h_{a,b}(z)+h_{a,b}(-z)}2
 =\frac{e^{-z^2/(4A)}}{2\sqrt{\pi A}}
 \left[\erfc\left(\frac{b|z|}{2a\sqrt A}\right)
       +\erfc\left(\frac{a|z|}{2b\sqrt A}\right)\right].
\end{equation}
For $c>0$ and $z\in\R$, one has the positive representation
\begin{equation}\label{eq:erfc}
 \erfc(c|z|)=\frac2\pi\int_0^\infty
          \frac{e^{-c^2z^2(1+v^2)}}{1+v^2}\,dv.
\end{equation}
Indeed, both sides equal one at $z=0$. For $z>0$, differentiating the
right side and evaluating the Gaussian integral gives
$-2c\pi^{-1/2}e^{-c^2z^2}$, the derivative of the left side.
Continuity at zero follows by dominated convergence.

Insert \eqref{eq:erfc} into \eqref{eq:density}. The two Gaussian exponents
are $-\kappa_1(v)z^2$ and $-\kappa_2(v)z^2$, where
\[
 \kappa_1(v)=\frac{A+b^2v^2}{4a^2A},\qquad
 \kappa_2(v)=\frac{A+a^2v^2}{4b^2A}.
\]
The resulting integrand is nonnegative and integrates to one over
$(z,v)$ because $k_{a,b}$ is a probability density. Consequently its
cosine transform may be interchanged with the $v$ integral by absolute
integrability. The Gaussian Fourier identity
\[
 \int_{\R}e^{-\kappa z^2}\cos(\sqrt t\,z)\,dz
 =\sqrt{\frac\pi\kappa}\,e^{-t/(4\kappa)}
\]
now gives \eqref{eq:kernel}.
Both weights there are positive. At $t=0$ their combined mass is one,
either by the probability-density calculation or by
\[
 \frac2\pi\int_0^\infty
 \frac{a\,dv}{(1+v^2)\sqrt{A+b^2v^2}}
 =\frac2\pi\arcsin\frac a{\sqrt A},
 \qquad
 \arcsin\frac a{\sqrt A}+\arcsin\frac b{\sqrt A}=\frac\pi2.
\]
Push these weights forward by the respective rates
$\lambda_1=a^2A/(A+b^2v^2)$ and
$\lambda_2=b^2A/(A+a^2v^2)$.
Every rate is positive for finite $v$, and each rate varies strictly on
$v>0$. Hence the resulting measure is a probability measure on
$(0,\infty)$ and is nondegenerate (in fact, atomless).
\end{proof}

\subsection*{Proof of the conjecture}
\begin{proof}[Proof of Theorem \ref{thm:main}]
Splitting the defining series into even and odd powers gives
\[
 E_\alpha(i\sqrt t)=E_{2\alpha,1}(-t)
                +i\sqrt t\,E_{2\alpha,1+\alpha}(-t).
\]
Both functions on the right have real values for $t\geq0$, so
$f_\alpha(t)=|E_\alpha(i\sqrt t)|^2$.
Take independent copies $(R,H)$ and $(R',H')$ from Lemma \ref{lem:scale}.
Since $X=RH$ and $X'=R'H'$ are identically distributed,
\[
 f_\alpha(t)=\E e^{i\sqrt t(X-X')}
            =\E\cos\bigl(\sqrt t(X-X')\bigr)
            =\E_{R,R'}K_{R,R'}(t).
\]
Here the real part is essential: the characteristic function conditional
on two unequal scales need not be real. The conditional cosine transform
is exactly the kernel of Lemma \ref{lem:kernel}.

Define $\nu_\alpha(B)=\E\nu_{R,R'}(B)$ for Borel sets
$B\subset(0,\infty)$. The explicit integral \eqref{eq:kernel} makes
measurability immediate. Tonelli's theorem gives \eqref{eq:laplace}.
All scales are finite and strictly positive almost surely. Each
conditional measure is atomless, so $\nu_\alpha$ is also atomless,
in particular nondegenerate.
For every $t>0$, differentiation under the integral is justified in a
neighborhood of $t$ by boundedness of $\lambda^n e^{-t\lambda/2}$.
Thus
\[
 (-1)^n f_\alpha^{(n)}(t)
  =\int_0^\infty\lambda^n e^{-t\lambda}\,\nu_\alpha(d\lambda)>0.
\]
Finally,
\[
 f_\alpha f_\alpha''-(f_\alpha')^2
 =\frac12\iint(\lambda-\mu)^2e^{-t(\lambda+\mu)}
       \,\nu_\alpha(d\lambda)\nu_\alpha(d\mu)>0,
\]
because $\nu_\alpha$ is not concentrated at a point. This proves strict
log-convexity as well.
\end{proof}

At the endpoint $\alpha=1/2$, the scales are $R=R'=1$. In this case
\eqref{eq:kernel} gives the particularly simple probability density
\begin{equation}\label{eq:endpoint}
 f_{1/2}(t)=\int_0^1
     e^{-t\lambda}\frac{2\,d\lambda}{\pi(2-\lambda)\sqrt{1-\lambda}}.
\end{equation}
To obtain it, set $a=b=1$ and $\lambda=2/(2+v^2)$ in
\eqref{eq:kernel}. Its mass is one, as is also seen by
$\lambda=1-u^2$, which yields $4\pi^{-1}\int_0^1(1+u^2)^{-1}du=1$.

\subsection*{The proposed spectral logarithmic convexity estimate}
\begin{corollary}
Let $(e_j)$ be an orthonormal basis of a Hilbert space, let $\mu_j\in\R$,
and let $\sum_j|c_j|^2<\infty$. For $0<\alpha\leq1/2$, define
\[
 u(t)=\sum_j c_jE_\alpha(-i\mu_jt^\alpha)e_j,\qquad t\geq0.
\]
Then for every $T>0$,
\[
 \|u(t)\|\leq\|u(0)\|^{1-t/T}\|u(T)\|^{t/T}
 \qquad(0\leq t\leq T).
\]
In particular this applies to the spectral solution in the source's
proposed estimate (5.1).
\end{corollary}
\begin{proof}
Theorem \ref{thm:main} implies $0<f_\alpha(s)\leq1$, so the series
converges in the Hilbert space and
\[
 \|u(t)\|^2=\sum_j|c_j|^2 f_\alpha(\mu_j^2t^{2\alpha}).
\]
The function $h=\log f_\alpha$ is decreasing and convex on $[0,\infty)$,
using continuity at zero. Since $t\mapsto t^{2\alpha}$ is concave,
$t\mapsto h(\mu_j^2t^{2\alpha})$ is convex: apply concavity first,
then monotonicity of $h$, and finally convexity of $h$.
With $\theta=t/T$, this gives
$f_\alpha(\mu_j^2t^{2\alpha})
 \leq f_\alpha(\mu_j^2T^{2\alpha})^\theta$.
H\"older's inequality for the sum yields
\[
 \|u(t)\|^2\leq
 \left(\sum_j|c_j|^2\right)^{1-\theta}
 \left(\sum_j|c_j|^2f_\alpha(\mu_j^2T^{2\alpha})\right)^\theta.
\]
The endpoints and the zero initial datum are immediate; taking square
roots completes the proof.
\end{proof}

\section*{Verification and scope}
The proof treats the full stated parameter range, including $\alpha=1/2$,
and every derivative order. The accompanying script performs 133 checks
at 75-digit precision: direct convolution, conditional characteristic
functions versus the Laplace integral, derivatives through order four,
the endpoint density, an independent $\alpha=1/4$ scale integral, and
the gamma moment identities. These finite checks supplement the proof.
The dated literature audit records current open status and no located
existing solution claim; it does not assert exhaustive knowledge of
unindexed or private work. No external peer review is claimed.

\section*{Task 2: Determine the remaining parameter range}
Characterize all $\alpha\in(1/2,1)$ for which
$t\mapsto|E_\alpha(i\sqrt t)|^2$ is completely monotone or log-convex.
Give a rigorous obstruction when either property fails and account
separately for the constant case $\alpha=1$.

\section*{Task 3: Compute the general Laplace measure}
Derive a direct density or useful special-function expression for
$\nu_\alpha$ when $0<\alpha<1/2$. Determine its support and endpoint
asymptotics, recovering \eqref{eq:endpoint} as $\alpha\uparrow1/2$.

\section*{Task 4: Classify completely monotone fractional powers}
Determine the exponents $q>0$ for which $f_\alpha^q$ is completely
monotone. In particular settle $q=1/2$, and characterize whether the
probability measure $\nu_\alpha$ is infinitely divisible.

\section*{Task 5: Quantify spectral stability from the representation}
Use the Laplace measure to obtain quantitative refinements of the
spectral logarithmic convexity estimate under assumptions on the
spectral support of the initial datum. Determine equality cases and
the resulting stability constants for the backward problem.

\begin{thebibliography}{9}
\bibitem{source}
S. E. Chorfi, F. Et-tahri, L. Maniar, and M. Yamamoto,
\emph{Forward and backward problems for abstract time-fractional
Schr\"odinger equations}, \href{https://arxiv.org/abs/2510.03600}{arXiv:2510.03600v1}
(2025); published in \emph{Communications in Nonlinear Science and
Numerical Simulation} \textbf{157} (2026), 109758,
\href{https://doi.org/10.1016/j.cnsns.2026.109758}{doi:10.1016/j.cnsns.2026.109758}.
\bibitem{followup}
S. E. Chorfi, F. Et-tahri, L. Maniar, and M. Yamamoto,
\emph{Maximal regularity for time-fractional Schr\"odinger equations
and application to nonlinear equations},
\href{https://arxiv.org/abs/2603.16726}{arXiv:2603.16726v1} (2026).
\bibitem{janson}
S. Janson, \emph{Moments of Gamma type and the Brownian supremum process
area}, \emph{Probability Surveys} \textbf{7} (2010), 1--52,
\href{https://doi.org/10.1214/10-PS160}{doi:10.1214/10-PS160};
\href{https://arxiv.org/abs/1002.4135}{arXiv:1002.4135}.
\end{thebibliography}
\end{document}
